\chapter{Random Forest: complete case study}
\label{ch:rf}

\begin{saybox}[title=What to say at the start of this section]
Give the intuition for a single tree first --- it partitions the
formulation space into regions and predicts the average shelf life inside
each region --- then explain that a forest averages many such partitions
built on perturbed data. Only then bring in the notation. Use project
features in the examples: a split on storage temperature, a split on pH.
\end{saybox}

\section{How Random Forest works}

\subsection{A single regression tree}

A regression tree recursively partitions the feature space. At each node
it searches over features $j$ and thresholds $t$ for the split that most
reduces the sum of squared deviations within the resulting children:
\begin{equation}
(j^\star, t^\star) = \argmin_{j, t}
\left[
\sum_{i \in L(j,t)} (y_i - \bar{y}_{L})^2
+
\sum_{i \in R(j,t)} (y_i - \bar{y}_{R})^2
\right],
\end{equation}
where $L$ and $R$ are the rows falling left and right of the threshold.
For the data at hand a first split might be on
\texttt{storage\_temperature\_c} at $\approx 8^\circ$C, separating chilled
from abused storage; a subsequent split might be on
\texttt{matrix\_water\_activity} or on the one-hot indicator for
\texttt{indicator\_type = clostridia}. Prediction for a new row is the mean
target of the training rows in the leaf it falls into:
\begin{equation}
\hat{f}_{\text{tree}}(\mathbf{x}) = \frac{1}{|\{i : \mathbf{x}_i \in \ell(\mathbf{x})\}|}
\sum_{i \,:\, \mathbf{x}_i \in \ell(\mathbf{x})} y_i .
\end{equation}

Two properties of this rule matter for everything in \Cref{ch:folds}:
splitting is \emph{axis-aligned}, and prediction is a \emph{local average
of training targets}. The second property has a hard consequence:
\begin{equation}
\min_i y_i \;\le\; \hat{f}_{\text{tree}}(\mathbf{x}) \;\le\; \max_i y_i
\qquad \text{for all } \mathbf{x},
\label{eq:rf-bound}
\end{equation}
a tree, and therefore a forest, \textbf{can never predict outside the range
of the training targets}. It cannot extrapolate. Remember this equation;
it explains the single worst fold in the entire study.

\subsection{From tree to forest}

The forest draws $B$ bootstrap resamples of the training set, fits an
unpruned tree to each while restricting each split to a random subset of
features, and averages:
\begin{equation}
\hat{f}_{\text{RF}}(\mathbf{x}) = \frac{1}{B} \sum_{b=1}^{B} \hat{f}_{\text{tree}}^{(b)}(\mathbf{x}) .
\end{equation}
If the individual trees have variance $\sigma^2$ and pairwise correlation
$\rho$, the ensemble variance is
\begin{equation}
\operatorname{Var}\!\left[\hat{f}_{\text{RF}}\right]
= \rho \sigma^2 + \frac{1-\rho}{B}\sigma^2
\;\xrightarrow[B \to \infty]{}\; \rho\sigma^2 .
\label{eq:rf-variance}
\end{equation}
Feature subsampling exists to reduce $\rho$: it is the only term that
survives in the limit. This is why adding trees always helps a little and
never hurts, and why the production setting of $400$ trees is a
computational choice rather than a regularization choice.

Averaging over the bound in \cref{eq:rf-bound} preserves it: the ensemble
is also confined to the convex hull of the training targets.

\section{Production configuration, recovered from source}

From \evd{train\_specialists.py}:
\begin{center}\small
\texttt{RandomForestRegressor(n\_estimators=400, max\_depth=None,\\
min\_samples\_leaf=2, random\_state=42, n\_jobs=-1)}
\end{center}

\begin{description}[leftmargin=1.5em, style=nextline, font=\bfseries\color{dgreen}]
\item[\texttt{n\_estimators=400}] $B = 400$ in \cref{eq:rf-variance}; the
  $\frac{1-\rho}{B}\sigma^2$ term is essentially exhausted at this value.
\item[\texttt{max\_depth=None}] trees grow until leaves are pure or cannot
  be split. Combined with the next parameter this is the highest-capacity
  setting available.
\item[\texttt{min\_samples\_leaf=2}] leaves may contain as few as two
  training rows --- minimal smoothing, near-maximal flexibility.
\item[\texttt{random\_state=42}] shared with the splitter, so the whole
  pipeline is reproducible from one seed.
\end{description}

The leave-rule-out experiments use \emph{the same} configuration, verified
by reading the model factory rather than trusting its name, so the
comparison between protocols is not confounded by a hyperparameter change.
The only differences between the production run and the diagnostic run are
the partition assignment and the fact that the diagnostic harness refits
the preprocessing inside every fold.

\section{Original Random Forest results}

\begin{table}[htbp]\centering\small
\caption{Random Forest under the original context-holdout protocol, V7
dataset. Extracted from the stored production metrics.}
\label{tab:rf_original}
\begin{tabular}{lrrrrrr}
\toprule
Specialist & $n_{\text{train}}$ & $n_{\text{test}}$ & train $\rsq$ & val $\rsq$ & test $\rsq$ & test RMSE (d) \\
\midrule
soft / general & 6\,160 & 1\,320 & 0.997 & 0.955 & \textbf{0.956} & 5.99 \\
semi-hard / general & 5\,040 & 1\,080 & 0.995 & 0.925 & \textbf{0.905} & 14.57 \\
hard / general & 4\,200 & 900 & 0.994 & 0.910 & \textbf{0.927} & 11.68 \\
soft / safety & 3\,360 & 720 & 0.997 & 0.548 & \textbf{0.572} & 10.62 \\
semi-hard / safety & 2\,800 & 600 & 0.997 & 0.015 & \textbf{0.342} & 15.76 \\
hard / safety & 2\,240 & 480 & 0.996 & $-0.131$ & \textbf{$-0.097$} & 21.86 \\
\bottomrule
\end{tabular}
\end{table}

The training $\rsq$ of $0.994$--$0.997$ across every specialist is exactly
what \texttt{max\_depth=None, min\_samples\_leaf=2} predicts and carries no
information about generalization. The informative contrast is between the
quality specialists (test $\rsq \ge 0.905$) and the safety specialists
(test $\rsq$ from $0.572$ down to $-0.097$), under an identical protocol
and an identical model.

\section{Revised Random Forest results}

The same configuration was re-evaluated under leave-one-rule-out
cross-validation across all $56$ folds, with every out-of-fold prediction
retained.

\begin{table}[htbp]\centering\small
\caption{Random Forest: original context-holdout against leave-rule-out,
per specialist. Pooled OOF $\rsq$ scores all out-of-fold predictions
together; mean and median are over folds.}
\label{tab:rf_revised}
\begin{tabular}{lrrrrrr}
\toprule
Specialist & folds & orig.\ test $\rsq$ & pooled OOF $\rsq$ & mean fold $\rsq$ & median fold $\rsq$ & OOF MAE (d) \\
\midrule
soft / general & 7 & 0.956 & \textbf{0.954} & 0.604 & 0.875 & 3.23 \\
semi-hard / general & 12 & 0.905 & \textbf{0.921} & 0.737 & 0.782 & 8.81 \\
hard / general & 8 & 0.927 & \textbf{0.582} & $-0.097$ & 0.630 & 14.56 \\
soft / safety & 11 & 0.572 & \textbf{0.844} & 0.692 & 0.812 & 4.39 \\
semi-hard / safety & 8 & 0.342 & \textbf{0.755} & 0.812 & 0.856 & 7.07 \\
hard / safety & 10 & $-0.097$ & \textbf{0.642} & 0.743 & 0.790 & 9.13 \\
\bottomrule
\end{tabular}
\end{table}

\noindent Full per-fold mean/median/min/max values appear in
\Cref{tab:tab_loro_ensembles_by_specialist}; the per-fold detail is in
\Cref{tab:tab_folds_random_forest}.

\begin{findingbox}[title=Two results that must be reported together]
\textbf{(a) The quality specialists hold up better than expected.} Soft and
semi-hard general-shelf-life retain pooled out-of-fold $\rsq$ of $0.954$
and $0.921$ --- essentially unchanged from their context-holdout figures.
Withholding an entire generation rule barely hurts them.

\textbf{(b) Hard/general collapses, and the safety specialists improve.}
Hard/general falls from $0.927$ to $0.582$ pooled. Meanwhile all three
safety specialists \emph{rise} sharply --- hard/safety from $-0.097$ to
$+0.642$. The revised protocol is not uniformly harsher; it is differently
structured, and for the safety specialists the original single random
context split happened to produce a harder and less representative test
partition than the rule-wise folds do.
\end{findingbox}

The second half of that finding deserves emphasis at the meeting because
it cuts against the expected narrative. The honest statement is not
``stricter evaluation, lower numbers everywhere.'' It is that the original
single split gave an unrepresentative reading in \emph{both} directions,
and that averaging over $56$ structured folds is more informative than one
random partition in either case.

\fig{fig20_rf_fold_r2_by_rule}{0.98}{%
  Random Forest leave-rule-out $\rsq$ for every one of the $56$ held-out
  rules, grouped by specialist and clipped at $-2$. Most rules are
  predicted well; a small number fail badly, and \Cref{ch:folds}
  investigates each of them individually.}

\fig{fig24_avp_soft_general}{0.62}{%
  Random Forest out-of-fold predictions against actual shelf life for the
  soft/general specialist, coloured by held-out rule. All $8\,800$ rows
  appear exactly once, each predicted by a model that had never seen its
  rule.}

\fig{fig25_avp_hard_general}{0.62}{%
  The same plot for hard/general. The detached cluster in the upper region
  --- predicted far below the diagonal --- is a single held-out rule, and
  it is responsible for most of this specialist's loss.}

\fig{fig26_residual_histograms}{0.98}{%
  Out-of-fold residual distributions per specialist. The hard/general panel
  shows the pronounced negative tail produced by that same rule.}

\fig{fig27_residual_vs_actual}{0.98}{%
  Residuals against the true value for the three quality specialists. The
  fan pattern in hard/general --- errors growing with the true value --- is
  the signature of \cref{eq:rf-bound} binding.}

\section{Where the error concentrates}

\fig{fig28_error_by_product}{0.92}{%
  Mean absolute out-of-fold error by cheese product, restricted to
  products with at least $40$ out-of-fold rows. Bar colour encodes the
  direction of the average bias.}

\fig{fig29_error_by_ingredient}{0.92}{%
  Mean absolute out-of-fold error by primary treatment ingredient, same
  restriction.}

\fig{fig30_signed_error_by_rule}{0.78}{%
  Mean signed error for every held-out rule. Systematic under-prediction
  (burgundy) dominates precisely where the held-out rule's targets lie
  above the training range --- again \cref{eq:rf-bound}.}
